\documentclass[10pt,a4paper]{amsart}
\usepackage[margin=19mm]{geometry}
\usepackage{amsmath,amssymb,mathtools}
\usepackage[scr=rsfs,cal=euler]{mathalfa}
\usepackage{xfrac}
\usepackage[unicode,hidelinks,bookmarks=false]{hyperref}
\newcommand{\ie}{i.e.\ }
\newcommand{\cf}{cf.\ }
\newcommand{\resp}{respectively\ }
\newcommand{\Subsection}{\subsection}
\providecommand{\nobreakdash}{\nobreak\hbox{-}}
\setlength{\emergencystretch}{2em}
\multlinegap0pt
\usepackage{amssymb}
\newtheorem{corollary}{Corollary}
\title{On the number of missing integers in partitions}
\author{Subhash Chand Bhoria, Pramod Eyyunni, Subhrangsu Santra}
\date{Source: arXiv:2604.12557v1 (CC BY 4.0)}
\begin{document}
\maketitle

\noindent\emph{Source and licence.} On the number of missing integers in partitions. Version arXiv:2604.12557v1 (CC BY 4.0), distributed by arXiv under the Creative Commons Attribution 4.0 International licence, http://creativecommons.org/licenses/by/4.0/, as recorded in the arXiv licence metadata for this exact version and shown on the abstract page https://arxiv.org/abs/2604.12557v1. Full licence notice: \texttt{licenses/arxiv-2604.12557-CC-BY-4.0.txt}.\par\smallskip

\noindent\textit{Editorial scope.} Counted unit: the article's corollary overponemissing (source0.tex lines 438--444). The article's complete original proof of it is delivered with it (source0.tex lines 514--532, one \texttt{proof} environment of 19 lines, which the article places away from the statement). No mathematics is written, repaired or re-derived: the statement and the proof are the article's own lines, in the article's own notation, and no retained line is substituted or re-flowed.  Retained uncounted as the source-local context the proof needs: lines 424--426.  Nothing was omitted from any retained span, so the package carries no omission record.  No retained line contains a citation, so the wrapper carries no bibliography and no bibliography entry.  No retained line contains a figure environment, an image-inclusion command or drawing code, so no image is delivered and the package declares no asset dependency.  Editorial formatting only, in addition: an AMS \texttt{amsart} preamble that restates the article's own declarations and macros used by the retained lines, namely the environments \texttt{corollary} on separate counters, together with the standard packages those lines need.  The retained environments print in this document as Corollary 1 in order of appearance, whereas the article prints the same items with the numbering of its own full document.  The complete native source of the article as distributed in the arXiv e-print for this exact version is bundled unchanged as \texttt{sources/198411/source0.tex}.
\begin{align}
\sum_{n=0}^{\infty}\sum_{m=0}^{\infty} \overline{P}(n,m)w^mq^n=\frac{\big((w-2)q;q\big)_{\infty}}{(wq;q)_{\infty}}.\label{W_GFN_OP}
\end{align}

\begin{corollary}\label{overponemissing}
We have
\begin{align*}
\sum_{n=0}^{\infty}\overline{P}(n,1)q^n=(-2q;q)_{\infty}\sum_{k=1}^{\infty}\frac{2q^{2k}}{1+2q^k}.
\end{align*}
Here $\overline{P}(n,1)$ is the number of overpartitions of $n$ in which exactly one integer is missing. This quantity is equinumerous with the number of overpartitions in which exactly one integer occurs twice and every other integer occurs only once, either in the non-overlined or overlined versions, but not both.
\end{corollary}

\begin{proof}[Corollary \textup{\ref{overponemissing}}][]
Differentiate both sides of \eqref{W_GFN_OP} with respect to $w$ to obtain
\begin{align*}
\frac{d}{dw} \frac{((w-2)q; q)_{\infty}}{(wq; q)_{\infty}}&= \frac{1}{(wq; q)_{\infty}} \times ((w-2)q; q)_{\infty} \sum_{n=1}^{\infty}\frac{-q^n}{1 - (w-2) q^n} \\
&+ ((w-2)q; q)_{\infty} \times \frac{-1}{(wq; q)_{\infty}^2} \times (wq; q)_{\infty} \times \sum_{n=1}^{\infty} \frac{-q^n}{1 - wq^n} \\
&= -\frac{((w-2)q; q)_{\infty}}{(wq; q)_{\infty}}\sum_{n=1}^{\infty}\frac{q^n}{1 - (w-2) q^n} + \frac{((w-2)q; q)_{\infty}}{(wq; q)_{\infty}} \sum_{n=1}^{\infty} \frac{q^n}{1 - wq^n}.
\end{align*}
Now, we set $w=0$ and get
\begin{align*}
&-(-2q; q)_{\infty}\sum_{n=1}^{\infty} \frac{q^n}{1 + 2q^n} + (-2q; q)_{\infty}\sum_{n=1}^{\infty}q^n \\
%& =(-2q; q)_{\infty}\left(\sum_{n=1}^{\infty} -\frac{q^n}{1 + 2q^n} + q^n \right) \\
&  = (-2q; q)_{\infty}\sum_{n=1}^{\infty}\frac{2q^{2n}}{1 + 2q^n}.
\end{align*}
For the combinatorial interpretation stated in the corollary, we observe that
\begin{equation*}
(-2q; q)_{\infty}\sum_{n=1}^{\infty}\frac{2q^{2n}}{1 + 2q^n} = \sum_{n=1}^{\infty} 2q^{2n} \prod_{\substack{k=1 \\ k \neq n}}^{\infty} (1+2q^k).
\end{equation*}
The general term on the right hand side above, $2q^{2n} \prod_{\substack{k=1 \\ k \neq n}}^{\infty} (1+2q^k)$, generates the overpartitions where $n$ occurs exactly twice (either both non-overlined or one overlined and non-overlined) and all other parts are distinct, occuring at most once either as an overlined or a non-overlined part but not both. Thus, the sum over all positive integers $n$ generates the required type of overpartitions.
\end{proof}

\end{document}
